\pagestyle{empty}
\tight
\begin{tblr}{width=\textwidth,colspec={|Q[c,m,1.9cm]|X[l,m]|},hlines,vlines,hspan=minimal,row{1}={bg=headblue},rowsep=1pt,colsep=3pt}
\SetCell{c,m}\hfil\logo\hfil & \textbf{POLITEKNIK NEGERI MALANG}\par\textbf{JURUSAN TEKNOLOGI INFORMASI}\par\textbf{PROGRAM STUDI: D4 TEKNIK INFORMATIKA} \\
\SetCell[c=2]{c} \textbf{RENCANA PEMBELAJARAN SEMESTER (RPS)} & \\
\end{tblr}
\begin{longtblr}{width=\textwidth,colspec={|Q[l,m,1.9cm]|X[0.85,l,m]|X[1.45,l,m]|X[0.75,l,m]|X[1.15,l,m]|},hlines,vlines,hspan=minimal,rowsep=1pt,colsep=2pt,row{1,3}={bg=headgray,font=\bfseries}}
MATA KULIAH & KODE & BOBOT (SKS)/jam & SEMESTER & TGL. PENYUSUNAN \\
Praktikum Pemrograman Berbasis Objek & RTI253008 & 2 SKS/4 jam & 3 & 1 Juli 2025 \\
OTORISASI & Dosen Pengembang RPS & Koordinator RMK & \SetCell[c=2]{l} Ka PRODI &  \\
 & Agung Nugroho Pramudhita, S.T., M.T. &  & \SetCell[c=2]{l}{\vspace{8mm}\par Dr. Ely Setyo Astuti, S.T., M.T.} &  \\
\SetCell[r=12]{l} Capaian Pembelajaran (CP) & \SetCell[c=4]{l,bg=headgray,font=\bfseries} Capaian Pembelajaran Lulusan Program Studi (CPL-Prodi) &  &  &  \\
 & CPL02 & \SetCell[c=3]{l} Mampu menerapkan prinsip rekayasa sistem dan perangkat lunak untuk menghasilkan solusi aplikasi multi-platform sesuai kebutuhan. &  &  \\
 & CPL07 & \SetCell[c=3]{l} Mampu mengembangkan sistem komputasi cerdas dan melakukan analisis data berbasis kecerdasan buatan untuk pengambilan keputusan. &  &  \\
 & CPL09 & \SetCell[c=3]{l} Mampu menjelaskan cara kerja sistem komputer dan menerapkan pendekatan komputasi untuk menyelesaikan masalah. &  &  \\
 & \SetCell[c=4]{l,bg=headgray,font=\bfseries} Capaian Pembelajaran mata kuliah (CPL-MK) &  &  &  \\
 & CPMK0210 & \SetCell[c=3]{l} Mahasiswa mampu menerapkan prinsip fundamental OOP dalam implementasi kode Java yang terstruktur. &  &  \\
 & CPMK0704 & \SetCell[c=3]{l} Mahasiswa mampu menggunakan bahasa Java dengan berbagai paradigma OOP untuk menyelesaikan masalah komputasi berbasis kasus. &  &  \\
 & CPMK0903 & \SetCell[c=3]{l} Mahasiswa mampu membangun aplikasi Java lengkap berbasis GUI dan koneksi database menggunakan pendekatan OOP. &  &  \\
 & \SetCell[c=4]{l,bg=headgray,font=\bfseries} Kemampuan akhir tiap tahapan belajar (Sub-CPMK) &  &  &  \\
 & SCPMK0210-02601 & \SetCell[c=3]{l} Mahasiswa mampu mengimplementasikan kelas, enkapsulasi, relasi kelas, dan inheritance dalam kode Java yang berfungsi. &  &  \\
 & SCPMK0704-02602 & \SetCell[c=3]{l} Mahasiswa mampu mengimplementasikan overriding, overloading, interface, kelas abstrak, dan polymorfisme pada kasus nyata. &  &  \\
 & SCPMK0903-02603 & \SetCell[c=3]{l} Mahasiswa mampu membangun aplikasi Java berbasis GUI (Swing/JavaFX) dengan koneksi database JDBC. &  &  \\
\SetCell[r=2]{l} Penilaian & \SetCell[c=4]{l} *Nilai CPMK didapat dari agregasi nilai SUB-CPMK yang dibebankan &  &  &  \\
 & \SetCell[c=4]{l}{\tiny\begin{tblr}{width=\linewidth,colspec={|Q[c,m,0.1\linewidth]|Q[c,m,0.16\linewidth]|X[l,m]|X[l,m]|X[l,m]|X[l,m]|X[l,m]|X[l,m]|Q[c,m,0.08\linewidth]|},hlines,vlines,hspan=minimal,rowsep=1pt,colsep=0.7pt,row{1,2}={bg=headgray,font=\bfseries}}
Kode CPMK & Kode SCPMK & \SetCell[c=6]{c} Bobot per Bentuk Penilaian & & & & & & Total Bobot \\
 & & Tugas & Kuis 1 & UTS & Kuis 2 & PBL & UAS & \\
CPMK0210 & SCPMK0210-02601 & 12\%: jobsheet OOP dasar & 5\%: konsep OOP dasar & 8\%: UTS OOP dasar & & & & 25\% \\
CPMK0704 & SCPMK0704-02602 & 10\%: jobsheet OOP lanjut & & 7\%: UTS OOP lanjut & 5\%: kuis OOP lanjut & 10\%: PBL GUI+database & & 32\% \\
CPMK0903 & SCPMK0903-02603 & 8\%: jobsheet GUI dan JDBC & & & & 12\%: PBL GUI+database & 23\%: demo aplikasi final & 43\% \\
\SetCell[c=8]{r}\textbf{Total Per Penilaian} & & & & & & & & \textbf{100\%} \\
\end{tblr}} &  &  &  \\
Diskripsi Singkat Mata Kuliah & \SetCell[c=4]{l} Praktikum Pemrograman Berbasis Objek memberikan pengalaman hands-on mengimplementasikan konsep OOP menggunakan Java, mencakup enkapsulasi, inheritance, polymorfisme, GUI Swing, dan koneksi database JDBC dalam proyek akhir terintegrasi. &  &  &  \\
Bahan Kajian / Materi Pembelajaran & \SetCell[c=4]{l}\begin{enumerate}\item Kelas, objek, enkapsulasi, konstruktor, dan relasi kelas\item Inheritance, overriding, overloading, kelas abstrak, dan interface\item Polymorfisme, exception handling, collections, dan generics\item GUI Swing: komponen dasar dan event handling\item Koneksi database JDBC dan aplikasi terintegrasi\end{enumerate} &  &  &  \\
\SetCell[r=2]{l} Pustaka & \SetCell[c=4]{l} \textbf{Utama:}\begin{enumerate}\item Deitel, P. \& Deitel, H. 2017. Java How to Program. Pearson.\item Bloch, J. 2018. Effective Java. Addison-Wesley.\item Dokumentasi Oracle Java SE dan JDBC.\end{enumerate} &  &  &  \\
 & \SetCell[c=4]{l} \textbf{Pendukung:} Materi LMS, dokumentasi resmi Java API, dan jobsheet praktikum. &  &  &  \\
Media Pembelajaran & \SetCell[c=2]{l} \textbf{Software:} JDK, IDE (Eclipse/IntelliJ/NetBeans), Scene Builder, LMS. &  & \SetCell[c=2]{l} \textbf{Hardware:} Komputer/laptop, proyektor. &  \\
Nama Dosen Pengampu & \SetCell[c=4]{l} Tim Pengajar Program Studi D4 TI &  &  &  \\
Matakuliah Syarat & \SetCell[c=4]{l} RTI251006 Dasar Pemrograman, RTI253007 Pemrograman Berbasis Objek &  &  &  \\
\end{longtblr}
\begin{longtblr}{width=\textwidth,colspec={|Q[c,m,0.06\textwidth]|X[1.35,l,m]|X[1.08,l,m]|X[1.16,l,m]|X[0.7,l,m]|X[1.24,l,m]|X[1.06,l,m]|X[0.98,l,m]|Q[c,m,0.05\textwidth]|},hlines,vlines,hspan=minimal,rowhead=2,row{1,2}={bg=headgreen,font=\bfseries},rowsep=1pt,colsep=1.5pt}
Mgg.\par Ke & Kemampuan Akhir Yang Direncanakan (Sub-CP-MK) & Bahan kajian (Materi Pembelajaran) & Bentuk dan Metode Pembelajaran & Estimasi Waktu & Pengalaman Belajar Mahasiswa & Kriteria \& Bentuk Penilaian & Indikator Penilaian & Bobot Penilaian (\%) \\
(1) & (2) & (3) & (4) & (5) & (6) & (7) & (8) & (9) \\
1 & SCPMK0210-02601 & Pengenalan lingkungan praktikum Java dan IDE. & Modalitas: Blended Learning. Bentuk: Luring/Daring. Metode: Case Method, praktik pemrograman Java, studi kasus, dan peer review. & 1 x 4 x 50' praktikum; 1 x 1 x 50' tugas/laporan mandiri. & Mahasiswa mengerjakan jobsheet pengenalan IDE dan membuat program Java pertama, menunjukkan evidence hasil belajar. & Bentuk penilaian: Tugas Jobsheet. Mengacu rubrik RTM. & Ketepatan konsep; kualitas implementasi; validasi dan dokumentasi. & 3 \\
2 & SCPMK0210-02601 & Praktikum: kelas dan objek. & Modalitas: Blended Learning. Bentuk: Luring/Daring. Metode: Case Method, praktik pemrograman Java, studi kasus, dan peer review. & 1 x 4 x 50' praktikum; 1 x 1 x 50' tugas/laporan mandiri. & Mahasiswa mengerjakan jobsheet kelas dan objek serta menunjukkan evidence hasil belajar. & Bentuk penilaian: Tugas Jobsheet. Mengacu rubrik RTM. & Ketepatan konsep; kualitas implementasi; validasi dan dokumentasi. & 3 \\
3 & SCPMK0210-02601 & Praktikum: enkapsulasi dan konstruktor. & Modalitas: Blended Learning. Bentuk: Luring/Daring. Metode: Case Method, praktik pemrograman Java, studi kasus, dan peer review. & 1 x 4 x 50' praktikum; 1 x 1 x 50' tugas/laporan mandiri. & Mahasiswa mengerjakan jobsheet enkapsulasi dan konstruktor serta menunjukkan evidence hasil belajar. & Bentuk penilaian: Tugas Jobsheet. Mengacu rubrik RTM. & Ketepatan konsep; kualitas implementasi; validasi dan dokumentasi. & 3 \\
4 & SCPMK0210-02601 & Praktikum: relasi kelas dan array objek. & Modalitas: Blended Learning. Bentuk: Luring/Daring. Metode: Case Method, praktik pemrograman Java, studi kasus, dan peer review. & 1 x 4 x 50' praktikum; 1 x 1 x 50' tugas/laporan mandiri. & Mahasiswa mengerjakan jobsheet relasi kelas dan array objek serta menunjukkan evidence hasil belajar. & Bentuk penilaian: Tugas Jobsheet. Mengacu rubrik RTM. & Ketepatan konsep; kualitas implementasi; validasi dan dokumentasi. & 3 \\
5 & SCPMK0210-02601 & Kuis 1 OOP dasar. & Modalitas: Blended Learning. Bentuk: Luring/Daring. Metode: Case Method, praktik pemrograman Java, studi kasus, dan peer review. & 1 x 4 x 50' praktikum; 1 x 1 x 50' tugas/laporan mandiri. & Mahasiswa mengerjakan kuis untuk mengukur pemahaman konsep OOP dasar. & Bentuk penilaian: Kuis 1. Mengacu rubrik RTM. & Ketepatan konsep; kualitas implementasi; validasi dan dokumentasi. & 5 \\
6 & SCPMK0704-02602 & Praktikum: inheritance dan overriding. & Modalitas: Blended Learning. Bentuk: Luring/Daring. Metode: Case Method, praktik pemrograman Java, studi kasus, dan peer review. & 1 x 4 x 50' praktikum; 1 x 1 x 50' tugas/laporan mandiri. & Mahasiswa mengerjakan jobsheet inheritance dan overriding serta menunjukkan evidence hasil belajar. & Bentuk penilaian: Tugas Jobsheet. Mengacu rubrik RTM. & Ketepatan konsep; kualitas implementasi; validasi dan dokumentasi. & 2 \\
7 & SCPMK0704-02602 & Praktikum: kelas abstrak dan interface. & Modalitas: Blended Learning. Bentuk: Luring/Daring. Metode: Case Method, praktik pemrograman Java, studi kasus, dan peer review. & 1 x 4 x 50' praktikum; 1 x 1 x 50' tugas/laporan mandiri. & Mahasiswa mengerjakan jobsheet kelas abstrak dan interface serta menunjukkan evidence hasil belajar. & Bentuk penilaian: Tugas Jobsheet. Mengacu rubrik RTM. & Ketepatan konsep; kualitas implementasi; validasi dan dokumentasi. & 2 \\
8 & SCPMK0210-02601, SCPMK0704-02602 & UTS Praktikum OOP dasar. & Modalitas: Blended Learning. Bentuk: Luring/Daring. Metode: Case Method, praktik pemrograman Java, studi kasus, dan peer review. & 1 x 4 x 50' praktikum; 1 x 1 x 50' tugas/laporan mandiri. & Mahasiswa mengerjakan ujian tengah semester berbasis praktikum untuk mengukur kemampuan OOP dasar. & Bentuk penilaian: UTS Praktikum. Mengacu rubrik RTM. & Ketepatan konsep; kualitas implementasi; validasi dan dokumentasi. & 15 \\
9 & SCPMK0704-02602 & Praktikum: polymorfisme. & Modalitas: Blended Learning. Bentuk: Luring/Daring. Metode: Case Method, praktik pemrograman Java, studi kasus, dan peer review. & 1 x 4 x 50' praktikum; 1 x 1 x 50' tugas/laporan mandiri. & Mahasiswa mengerjakan jobsheet polymorfisme serta menunjukkan evidence hasil belajar. & Bentuk penilaian: Tugas Jobsheet. Mengacu rubrik RTM. & Ketepatan konsep; kualitas implementasi; validasi dan dokumentasi. & 2 \\
10 & SCPMK0704-02602 & Praktikum: exception handling. & Modalitas: Blended Learning. Bentuk: Luring/Daring. Metode: Case Method, praktik pemrograman Java, studi kasus, dan peer review. & 1 x 4 x 50' praktikum; 1 x 1 x 50' tugas/laporan mandiri. & Mahasiswa mengerjakan jobsheet exception handling serta menunjukkan evidence hasil belajar. & Bentuk penilaian: Tugas Jobsheet. Mengacu rubrik RTM. & Ketepatan konsep; kualitas implementasi; validasi dan dokumentasi. & 2 \\
11 & SCPMK0704-02602 & Praktikum: collections dan generics. & Modalitas: Blended Learning. Bentuk: Luring/Daring. Metode: Case Method, praktik pemrograman Java, studi kasus, dan peer review. & 1 x 4 x 50' praktikum; 1 x 1 x 50' tugas/laporan mandiri. & Mahasiswa mengerjakan jobsheet collections dan generics serta menunjukkan evidence hasil belajar. & Bentuk penilaian: Tugas Jobsheet. Mengacu rubrik RTM. & Ketepatan konsep; kualitas implementasi; validasi dan dokumentasi. & 2 \\
12 & SCPMK0704-02602 & Kuis 2 OOP lanjut. & Modalitas: Blended Learning. Bentuk: Luring/Daring. Metode: Case Method, praktik pemrograman Java, studi kasus, dan peer review. & 1 x 4 x 50' praktikum; 1 x 1 x 50' tugas/laporan mandiri. & Mahasiswa mengerjakan kuis untuk mengukur pemahaman konsep OOP lanjut. & Bentuk penilaian: Kuis 2. Mengacu rubrik RTM. & Ketepatan konsep; kualitas implementasi; validasi dan dokumentasi. & 5 \\
13 & SCPMK0903-02603 & Praktikum: GUI Swing komponen dasar. & Modalitas: Blended Learning. Bentuk: Luring/Daring. Metode: Case Method, praktik pemrograman Java, studi kasus, dan peer review. & 1 x 4 x 50' praktikum; 1 x 1 x 50' tugas/laporan mandiri. & Mahasiswa mengerjakan jobsheet komponen GUI Swing dasar serta menunjukkan evidence hasil belajar. & Bentuk penilaian: Tugas Jobsheet. Mengacu rubrik RTM. & Ketepatan konsep; kualitas implementasi; validasi dan dokumentasi. & 3 \\
14 & SCPMK0903-02603 & Praktikum: GUI Swing event handling. & Modalitas: Blended Learning. Bentuk: Luring/Daring. Metode: Case Method, praktik pemrograman Java, studi kasus, dan peer review. & 1 x 4 x 50' praktikum; 1 x 1 x 50' tugas/laporan mandiri. & Mahasiswa mengerjakan jobsheet event handling pada GUI Swing serta menunjukkan evidence hasil belajar. & Bentuk penilaian: Tugas Jobsheet. Mengacu rubrik RTM. & Ketepatan konsep; kualitas implementasi; validasi dan dokumentasi. & 3 \\
15 & SCPMK0903-02603 & Praktikum: koneksi database JDBC. & Modalitas: Blended Learning. Bentuk: Luring/Daring. Metode: Case Method, praktik pemrograman Java, studi kasus, dan peer review. & 1 x 4 x 50' praktikum; 1 x 1 x 50' tugas/laporan mandiri. & Mahasiswa mengerjakan jobsheet koneksi database menggunakan JDBC serta menunjukkan evidence hasil belajar. & Bentuk penilaian: Tugas Jobsheet. Mengacu rubrik RTM. & Ketepatan konsep; kualitas implementasi; validasi dan dokumentasi. & 2 \\
16 & SCPMK0704-02602, SCPMK0903-02603 & PBL: Pengembangan aplikasi Java GUI+Database. & Modalitas: Blended Learning. Bentuk: Luring/Daring. Metode: Case Method, praktik pemrograman Java, studi kasus, dan peer review. & 1 x 4 x 50' praktikum; 1 x 1 x 50' tugas/laporan mandiri. & Mahasiswa mengembangkan aplikasi Java berbasis GUI dan database secara terintegrasi sebagai proyek akhir. & Bentuk penilaian: PBL Proyek Akhir. Mengacu rubrik RTM. & Ketepatan konsep; kualitas implementasi; validasi dan dokumentasi. & 22 \\
17 & SCPMK0903-02603 & UAS Praktikum: demo aplikasi final. & Modalitas: Blended Learning. Bentuk: Luring/Daring. Metode: Case Method, praktik pemrograman Java, studi kasus, dan peer review. & 1 x 4 x 50' praktikum; 1 x 1 x 50' tugas/laporan mandiri. & Mahasiswa mendemonstrasikan aplikasi Java final yang telah dikembangkan dan mempertahankan keputusan implementasi. & Bentuk penilaian: UAS Praktikum. Mengacu rubrik RTM. & Ketepatan konsep; kualitas implementasi; validasi dan dokumentasi. & 23 \\
\end{longtblr}
